\begin{quadro}[h!]
    \centering
    \caption{Exemplo de especificação textual de um Caso de Uso para a funcionalidade de saque}
    \label{qua:exemplo_caso_de_uso}
    
    \fcolorbox{black}{backcolour}{% 
        \begin{minipage}{0.9\textwidth}\small%
            \vspace{3pt}%
            \noindent\textbf{\normalsize CdU1 - Efetuar saque em caixa eletrônico}\vspace{5pt}\\
            \textbf{Agente:} Cliente\\
            \textbf{Pré-Condições:}
            \begin{itemize}[topsep=0pt]
                \itemsep0em
                \item O terminal do caixa eletrônico deve estar operante, conectado à rede bancária e abastecido com cédulas. 
                \item O cliente deve possuir um cartão válido e conhecer sua respectiva credencial de acesso (senha).
            \end{itemize}
            \textbf{Fluxo Principal:}
            \begin{enumerate}[topsep=0pt]
                \itemsep0em
                \item O cliente insere o cartão no terminal do caixa eletrônico.
                \item O sistema faz a leitura do cartão e solicita e apresenta o menu de operações.
                \item O cliente seleciona a opção "Saque".
                \item O sistema solicita a inserção do valor desejado.
                \item O cliente informa o valor a ser sacado.
                \item O sistema verifica a disponibilidade do saldo na conta, o limite de saque diário do cliente e a quantidade de cédulas disponíveis no terminal.
                \item O sistema registra a transação e deduz o valor solicitado do saldo da conta do cliente.
                \item O sistema dispensa as cédulas correspondentes no compartimento de saída.
                \item O sistema ejeta o cartão e encerra a sessão.
            \end{enumerate}
            \textbf{Fluxo de Exceção (Saldo Insuficiente):}\\
            No passo 8, caso o cliente não possua saldo ou limite diário suficiente, o sistema exibe uma mensagem informativa e ejeta o cartão.
        \end{minipage}
    }

    \vspace{5pt}
    \hspace*{0.75cm}\raggedright{\small \textbf{Fonte:} Elaborado pelo autor.}
\end{quadro}